\documentclass[11pt]{article}
\usepackage[a4paper,margin=1.45cm]{geometry}
\usepackage{fontspec}
\usepackage{xeCJK}
\setmainfont{TeXGyreTermes}
\setCJKmainfont{Noto Serif CJK JP}
\setsansfont{TeXGyreHeros}
\setCJKsansfont{Noto Sans CJK JP}
\usepackage{amsmath,amssymb,mathtools,array,booktabs}
\usepackage[protrusion=true,expansion=false]{microtype}
\setlength{\parindent}{1.0em}
\setlength{\parskip}{0pt}
\allowdisplaybreaks
\setlength{\emergencystretch}{30em}
\sloppy
\hbadness=10000
\vbadness=10000
\hfuzz=10pt
\vfuzz=10pt
\raggedbottom
\providecommand{\mG}{\mathfrak{G}}
\providecommand{\mA}{\mathfrak{A}}
\providecommand{\mB}{\mathfrak{B}}
\providecommand{\mC}{\mathfrak{C}}
\providecommand{\mD}{\mathfrak{D}}
\providecommand{\mE}{\mathfrak{E}}
\providecommand{\mH}{\mathfrak{H}}
\providecommand{\mK}{\mathfrak{K}}
\providecommand{\mM}{\mathfrak{M}}
\providecommand{\mN}{\mathfrak{N}}
\providecommand{\mR}{\mathfrak{R}}
\providecommand{\mS}{\mathfrak{S}}
\providecommand{\mT}{\mathfrak{T}}
\providecommand{\mU}{\mathfrak{U}}
\providecommand{\mO}{\mathfrak{O}}
\providecommand{\mo}{\mathfrak{o}}
\providecommand{\ma}{\mathfrak{a}}
\providecommand{\mb}{\mathfrak{b}}
\providecommand{\mc}{\mathfrak{c}}
\providecommand{\md}{\mathfrak{d}}
\providecommand{\mf}{\mathfrak{f}}
\providecommand{\mh}{\mathfrak{h}}
\providecommand{\ml}{\mathfrak{l}}
\providecommand{\mn}{\mathfrak{n}}
\providecommand{\mr}{\mathfrak{r}}
\providecommand{\mt}{\mathfrak{t}}
\providecommand{\mx}{\mathfrak{x}}
\providecommand{\mz}{\mathfrak{z}}
\providecommand{\mZ}{\mathfrak{Z}}
\providecommand{\mL}{\mathfrak{L}}
\providecommand{\iso}{\simeq}
\providecommand{\ann}{\operatorname{Ann}}
\providecommand{\tuple}[1]{(#1)}
\providecommand{\mat}[1]{\begin{pmatrix}#1\end{pmatrix}}
\begin{document}


\setcounter{footnote}{0}
\section*{34. 超複素量と表現論}
\emph{続きおよび完結：第III章 §§15--19 および第IV章 §§20--26。}

\subsection*{第III章．加群論と表現論}

\subsection*{§ 15. 表現と表現加群}

$\mo$ を環、$K$ を単位元をもつ環とする。後の応用では $K$ は常に体である。ただし非可換の文脈では、これは非可換体、すなわち斜体として読むべきである。

$\mo$ の $K$ における $n$ 次の表現とは、準同型
\[
        \mo \sim \mathfrak D,
\]
であり、ここで $\mathfrak D$ は $K$ 上の $n$ 次行列からなる環である。

$\mo$ の $K$ に関する\emph{表現加群}とは、$\mo$ 上の左加群かつ $K$ 上の右加群である双加群 $\mM$、
\[
        \mo\mM\subseteq \mM,
        \qquad
        \mM K\subseteq \mM,
\]
であって、さらに有限個の一生成元 $K$-加群の直和
\[
        \mM=x_1K+\cdots+x_nK
\]
であり、$K$ の単位元が恒等作用素であるものをいう。

\emph{すべての表現加群は表現を与える。} $c\in\mo$ とし、
\[
        cx_k=\sum_i x_i\gamma_{ik},
\]
または行列形式で
\[
        c(x_1,\ldots,x_n)=(x_1,\ldots,x_n)C,
        \qquad C=(\gamma_{ik})
\]
と書く。すると行列 $C$ は $c$ の表現をなす。実際、
\[
        (b+c)x_k=\sum_i x_i(\beta_{ik}+\gamma_{ik}),
\]
かつ
\[
\begin{aligned}
        bcx_k
        &=b\sum_jx_j\gamma_{jk}
          =\sum_j b x_j\gamma_{jk}
          =\sum_j\sum_i x_i\beta_{ij}\gamma_{jk}       \\
        &=\sum_i x_i\biggl(\sum_j\beta_{ij}\gamma_{jk}\biggr),
\end{aligned}
\]
すなわち
\[
        bc(x_1,\ldots,x_n)=(x_1,\ldots,x_n)BC
\]
である。

逆に、$\mo$ の任意の表現 $\mathfrak D$ は一つの表現加群に属し、しかもその加群のある特定の基底に属する。すなわち、$\mM$ を $x_1,\ldots,x_n$ の形式的線形形式全体
\[
        y=\sum_i x_i\alpha_i,
        \qquad \alpha_i\in K
\]
と解する。すると $\mM$ は右 $K$-加群である。さらに、元 $c$ に行列 $C=(\gamma_{ik})$ が対応するとき、
\[
        cx_k=\sum_i x_i\gamma_{ik},
        \qquad
        c\biggl(\sum_k x_k\alpha_k\biggr)
        =\sum_i x_i\sum_k\gamma_{ik}\alpha_k
\]
と定める。準同型関係
\[
        c+d\sim C+D,
        \qquad
        cd\sim CD
\]
から
\[
        (c+d)x_i=cx_i+dx_i,
        \qquad
        cdx_i=c(dx_i)
\]
が従い、同じことは和 $\sum_i x_i\alpha_i$ についても成り立つ。残る双加群の法則
\[
        c(y+z)=cy+cz,
        \qquad
        c(y\alpha)=(cy)\alpha
\]
は自明である。したがって $\mM$ は表現加群であり、構成によりちょうど表現 $\mo\sim\mathfrak D$ に属する。

$(y_1,\ldots,y_n)$ を同じ表現加群の別の基底とする：
\[
        (y_1,\ldots,y_n)=(x_1,\ldots,x_n)P,
        \qquad
        (x_1,\ldots,x_n)=(y_1,\ldots,y_n)P^{-1}.
\]
$x$-基底に関して $b\sim B$ であれば、$y$-基底に関しては
\[
        b(y_1,\ldots,y_n)
        =b(x_1,\ldots,x_n)P
        =(x_1,\ldots,x_n)BP
        =(y_1,\ldots,y_n)P^{-1}BP
\]
である。表現 $b\sim B$ と $b\sim P^{-1}BP$ を\emph{同値}と呼び、同じ表現類に数える。任意の正則行列 $P$ は $\mM$ の一つの基底を別の基底に移すので、つぎが示された。

\emph{すべての表現加群は、一意に定まる一つの表現類を与える。}\textsuperscript{15a)}

{\footnotesize\noindent
15) 可換体に関する表現加群は、W. Krull, \emph{Theorie und Anwendung der verallgemeinerten Abelschen Gruppen}, Sitzungsberichte der Heidelberger Akademie, 1926, 第1論文、に初めて現れる。\par
15a) 校正時追記（1929年4月14日）。B. L. van der Waerden が私に知らせてくれたように、「線形変換」と「行列」という概念を分ければ、特別な基底の選択に依存しない不変的な結びつきを得ることができる。線形変換とは二つの線形形式加群の準同型であり、行列とは、特定の基底の選択に対するこの準同型の表現である。I. すべての表現加群は、乗法子領域 $\mo$ に、加群自身の中への線形変換からなる一意な準同型系を割り当てる。実際、§2 の終わりにより、双加群 $\mM$ の左乗法子は、右作用素（ここでは $K$ による乗法）に関する $\mM$ から自身への作用素準同型を生成し、この対応は準同型である。II. 線形形式加群の自身の中への線形変換の系で、$\mo$ に準同型であるものはすべて、表現加群を与える。\par}

二つの表現加群が作用素同型であれば、それらが同じ表現類に対応することは明らかである。逆に、二つの表現加群 $(x_1,\ldots,x_n)$ と $(\bar x_1,\ldots,\bar x_n)$ が同じ表現を生成するなら、対応
\[
        x_i\mapsto \bar x_i,
        \qquad
        \sum_i x_i\alpha_i\mapsto \sum_i \bar x_i\alpha_i
\]
は同型を与える。乗法規則は行列 $C$ によって完全に決定されるからである。

特殊な場合として、$\mM$ の二つの基底が同じ表現を生成するなら、それらは $\mM$ から自身への作用素同型によって互いに移される。

同値な言い方をすれば、すべての $C$ について一つの固定された $P$ により $C=PCP^{-1}$ であるなら、対応
\[
        y_i\mapsto x_i,
        \qquad
        y_i=\sum_k x_k\pi_{ik}
\]
は $\mM$ から自身への同型である。逆に、双加群 $\mM$ の自身への任意の同型 $y_i\mapsto x_i$ は、すべての $C$ と可換な行列 $P$ によって媒介される。

表現の概念は、可換な乗法子領域 $P$ とまだ可換的に結びついている環 (§8) にも拡張できる。その場合、表現環 $K$ としては、その中心に $P$ を含むものだけを取り、表現には単に環準同型 $\mo\sim\mathfrak D$ であることだけではなく、作用素準同型であることを要求する。すなわち $a\sim A$ から
\[
        a\rho\sim A\rho\qquad(\rho\in P)
\]
が従うべきである。表現加群に対してこの要求は、追加の規則
\[
        am\cdot\rho=a(m\rho)=a\rho\cdot m
\]
を意味する。このとき加群と環は $P$ と\emph{可換的に結びついている}という。

\subsection*{§ 16. 可約表現}

$\mU$ が表現加群 $\mM$ の部分加群であり、$\mU$ の基底 $z_1,\ldots,z_t$ に $y_1,\ldots,y_r$ を付け加えて $\mM$ の基底を選べるとする。すなわち
\[
        \mM=y_1K+\cdots+y_rK+z_1K+\cdots+z_tK,        \tag{2a}
\]
である。このとき表現は
\[
        C=\begin{pmatrix}R&0\\ S&T\end{pmatrix}
\]
の形をもつ。ここで行列 $T$ だけで、$\mU$ によって媒介される $t$ 次の表現をなし、行列 $R$ は $\mM/\mU$ によって媒介される $r$ 次の表現をなす。

実際、
\[
        (cy_1,\ldots,cy_r,cz_1,\ldots,cz_t)
        =(y_1,\ldots,y_r,z_1,\ldots,z_t)C.
\]
$cz_i$ は $\mU$ の元として $z_i$ だけで表されるので、$C$ の右上ブロックは零である。$C$ の他の部分を上に示した順に $R,S,T$ と名づけるなら、とくに
\[
        (cz_1,\ldots,cz_t)=(z_1,\ldots,z_t)T
\]
である。したがって $T$ は $\mU$ によって媒介される表現をなす。さらに
\[
        (cy_1,\ldots,cy_r)\equiv (y_1,\ldots,y_r)R \pmod{\mU},
\]
一方で $y_i$ は $\mU$ を法として線形独立な基底をなす。よって $R$ は $\mM/\mU$ によって生成される表現をなす。

逆に、
\[
        C=\begin{pmatrix}R&0\\ S&T\end{pmatrix}
\]
という「可約」表現が与えられ、$R$ と $T$ が正方行列であるなら、対応する表現加群において、各 $c$ と最後の $t$ 個の基底元 $z_1,\ldots,z_t$ との積はそれらの元だけで表される。したがって $\mU=(z_1,\ldots,z_t)$ は部分加群である。

系。 
\[
        \mM\supset \mU_1\supset \mU_2\supset\cdots\supset \mU_e=0
\]
を $\mM$ の組成列とし、各 $\mU_i$ は直前のものの $K$-加群としての直和因子であると仮定する。このとき基底
\[
        z_{11},\ldots,z_{1r_1},\quad
        z_{21},\ldots,z_{2r_2},\quad\ldots\quad
        z_{e1},\ldots,z_{er_e}
\]
を、$i>\nu$ の $z_{ik}$ が $\mU_\nu$ の基底をなすように選ぶことができる。$\mM$ によって媒介される表現はそのとき
\[
 C=\begin{pmatrix}
 R_{11}&0&\cdots&0\\
 R_{12}&R_{22}&\cdots&0\\
 \vdots&\vdots&\ddots&\vdots\\
 R_{1e}&R_{2e}&\cdots&R_{ee}
 \end{pmatrix},                                      \tag{2}
\]
の形をもち、$R_{\nu\nu}$ は組成因子 $\mU_{\nu-1}/\mU_\nu$ によって媒介される表現をなす。

個々の表現 $R_{\nu\nu}$ は既約である。なぜなら組成因子 $\mU_{\nu-1}/\mU_\nu$ が単純だからである。

$K$ が体であると仮定し、したがって任意の部分加群が直和因子であるなら、逆に、形 (2) で可能な限り簡約された任意の表現は組成列を与える。Jordan--Hölder の定理により、組成因子は作用素同型を除いて一意に定まる。したがって $R_{\nu\nu}$ は同値な表現および順序を除いて一意に定まる。\textsuperscript{16)}

{\footnotesize\noindent 16) 言い換えれば、$\mU$ が $K$-加群として直和因子であるということである。$K$ が体なら、この仮定は常に満たされる。\par}

$\mM=\mA+\mB$ が二つの表現加群 $\mA=(y_1,\ldots,y_r)$ と $\mB=(z_1,\ldots,z_s)$ の直和であるなら、基底 $(y_1,\ldots,y_r,z_1,\ldots,z_s)$ によって媒介される表現は
\[
        C=\begin{pmatrix}R&0\\0&S\end{pmatrix}
\]
の形をもつ。ここで $R$ は $\mA$ によって媒介される元 $c$ の表現、$S$ は $\mB$ によって媒介される表現を表す。逆も同様である。まず $\mM$ を直接分解不能な成分の直和として書き、それらの中に上のような組成列を引けば、適切な基底に対して表現は、外側の大きな箱が直接分解不能な成分に対応し、対角の小ブロックが組成因子に対応するブロック形をもつ。

ふたたび $K$ が体であるなら、このように可能な限り簡約された任意の表現は、逆に、直接分解不能な和因子とその内部の組成因子による加群の表現を与える。Krull と Otto Schmidt の定理により、直接分解不能な表現成分の類は順序を除いて一意に定まる。

加群 $\mM$ が完全可約であるなら、ブロック形のすべての箱は一つの既約表現だけからなり、この表現を完全可約と呼ぶ。

\subsection*{§ 17. 単位元をもつ環に対する表現加群の直和分解}

$\mo$ を単位元をもつ環とし、$\mM$ を $\mo$-加群とする。このとき
\[
        \mM=\mM_e+\mM_0
\]
であり、単位元は $\mM_e$ 上で恒等作用素、$\mM_0$ 上で零作用素である。$\mM$ がさらに右 $K$-加群であるなら、$\mM_e$ と $\mM_0$ も右 $K$-加群である。

証明。任意の $m\in\mM$ は
\[
        m=em+(m-em)
\]
と書ける。$em$ 全体を $\mM_e$、$m-em$ 全体を $\mM_0$ とする。これらは明らかに加法群であり、$\mM$ が右 $K$-加群であれば右 $K$-加群でもある。さらに
\[
        r(em)=erm
\]
なので $\mM_e$ は $\mo$-加群でもあり、
\[
        r(m-em)=rm-rem=rm-rm
\]
なので $\mM_0$ も $\mo$-加群である。$em$ については $e$ が恒等作用素であり、$m-em$ については零作用素である。したがって $\mM_e$ と $\mM_0$ の両方に属するものは零だけで、和 $\mM_e+\mM_0$ は直和である。

表現加群については、$\mM_0$ は各元に零を割り当てる自明表現を与える。和因子 $\mM_0$ を分離すれば、$\mM_e$ によって媒介される表現が残り、そこで単位元は単位行列に対応する。したがって以後はそのような表現に制限してよいし、そうする。

\[
        \mo=\ma_1+\cdots+\ma_s
\]
を両側イデアルの直和とし、
\[
        e=e_1+\cdots+e_s
\]
を対応する $e$ の分解とする。また $\mM$ を、$e$ が恒等作用素である $\mo$-加群とする。このとき
\[
        \mM=\ma_1\mM+\cdots+\ma_s\mM
\]
は直和である。

実際、明らかに $\mM=\mo\mM=(\ma_1\mM,\ldots,\ma_s\mM)$ である。$\mb_i=\ma_1+\cdots+\widehat{\ma_i}+\cdots+\ma_s$ とおく。すると $\mb_i$ は対応して両側イデアルであり、
\[
        \mM=(\ma_i\mM,\mb_i\mM)
\]
で、この和は直和である。なぜなら $e_i$ は $\ma_i\mM$ 上で恒等作用素、$\mb_i\mM$ 上で零作用素だからである。

この定理に基づき、通常は両側分解不能な環の表現に制限する。

\subsection*{§ 18. 完全可約環の加群論と表現論}

$\mo$ を完全可約かつ両側単純とし、$\mM$ を $\mo$ の単位元が恒等作用素である有限 $\mo$-加群とする。このとき $\mM$ は完全可約であり、その単純構成要素は単純左イデアル $\ml_i$ と作用素同型である。

証明。
\[
        \mo=\ml_1+\cdots+\ml_n,
        \qquad
        \mM=(\mo m_1,\ldots,\mo m_k)
\]
とする。したがって
\[
        \mM=(\ldots,\ml_i m_k,\ldots).                         \tag{4}
\]
零でない $\ml_i m_k$ は $\ml_i$ と同型である。写像 $a\mapsto am_k$ が作用素同型だからである。したがって $\ml_i m_k$ は単純である。(4) から、すでに先行する和に含まれている $\ml_i m_k$ を省けば、和は直和になる。

したがって、さらに $\mM$ が直接分解不能であるなら、$\mM$ は単純であり、ある $\ml_i$ と同型である。

この単純な $\mM$ の自己同型斜体を $\Delta$、$\ml_i$ のそれを $\Lambda$ とする。$\mM$ と $\ml_i$ の作用素同型により、環同型 $\Delta\simeq\Lambda$ が得られる。§1 により、$\mM$ と $\ml_i$ はそれぞれ $\Delta$ と $\Lambda$ を右側の領域とする双加群とみなすことができ、したがって表現加群でもある。それらによって媒介される表現は、環同型 $\Delta\simeq\Lambda$ の下で互いに移る。ゆえにこの表現類を研究するには、$\ml_i$ によって、その自己同型斜体において媒介される表現に制限してよい。

具体的には、$\ml_1$ の行列単位の基底を
\[
        \ml_1=(c_{11},\ldots,c_{n1})
\]
の形に置き、前の §14 の自己同型斜体 $\Lambda$ を $\mo$ の反対斜体によって実現する。前の考察では、右イデアルを左イデアルに置き換え、自己同型を右に書くべきである。
\[
        c=\sum_{i,k}c_{ik}\alpha_{ik}
\]
が $\mo$ の元であるなら、
\[
        (cc_{11},\ldots,cc_{n1})=(c_{11},\ldots,c_{n1})(\alpha_{ik}).
\]
したがって $\ml_1$ によって媒介される表現は行列 $(\alpha_{ik})$ で与えられる。$\ml_\nu=(c_{1\nu},\ldots,c_{n\nu})$ についても、まったく同じ表現を得る。

$\mo$ が両側単純な成分の直和であるなら、これらの成分は順に同型的に表現され、残りの成分は零として作用する。

\subsection*{§ 19. 加群および表現加群における単純組成因子}

$\mo$ を単位元をもつ環、$\mc$ をその根基とする。$\mo$ 上の任意の単純加群および任意の単純表現加群において、$\mc$ の元は零化作用をもつ。したがって $\mo/\mc$ に対する加群または表現加群が得られる。

実際、$\mc\mM$ はふたたび加群である（表現の場合には右 $K$-加群でもある）。$\mM$ は単純なので、$\mc\mM=\mM$ または $\mc\mM=0$ のいずれかである。前者の場合には任意の指数 $\rho$ について $\mc^\rho\mM=\mM$ となり、$\mc$ の冪零性に矛盾する。したがって $\mc\mM=0$ である。

このことから、任意の加群および表現加群の単純組成因子もまた $\mo/\mc$ の単純左イデアルによって生成される。

$P$ と可換的に結びついた表現加群については、これはとくに、§16 のように組成列によって任意の表現を簡約すると、零を除いて、$\mo/\mc$ の単純左イデアルによって媒介される対角表現だけが現れることを意味する。

これらの事実には、$P$ 上の表現にすでに課した仮定を超える $\mo$ の有限性条件は不要である。任意の表現は、その絶対乗法子領域の同型表現である。この絶対乗法子領域は、有限行列環の逆像として、それ自身が有限階数の $P$-加群である。許されるイデアルは定義により $P$-加群であるから、極大条件と極小条件がそれについて成り立つ。

$P$ が可換体であるなら、この絶対乗法子領域は $P$ 上の超複素系である。この仮定は、対角上に非零表現が現れる場合には自動的に満たされる。そのとき $P$ は単純左イデアルの自己同型斜体の部分体に同型であり、$\mo/\mc$ の部分体 $K$ による実現において、可換的結合はそれを $K$ の中心の中へ強制するからである。

\subsection*{第IV章．群と超複素系の表現}

\subsection*{§ 20. 超複素系の組み入れ}

ここから、可換体 $P$ に関する超複素系 $\mo$ を考える。約束により、$\mo$ のイデアルとして数えるのは $P$-加群だけである。したがって左イデアルおよび右イデアルについて極大条件と極小条件が満たされ、第II章の環論全体が適用される。

以下で超複素系 $\mo$ の表現というときは、常に体 $P$ における表現、より正確には、$a\sim A$ から
\[
        a\rho\sim A\rho\qquad(\rho\in P)
\]
が従うものを意味する。すなわち $P$ に関する加群準同型性である。表現加群ではこれは
\[
        a\rho\cdot m=a(m\rho)
\]
と読まれる。したがって §19 の結果が適用される。すべての既約表現は
\[
        \mo_0=\mo/\mc
\]
の単純左イデアルによって媒介される。ただし $\mc$ は根基である。したがって、分解
\[
        \mo_0=\ma_1+\cdots+\ma_s
\]
に現れる両側単純和因子 $\ma$ の数だけ、互いに非同値な既約表現がある。

そのような左イデアルによって定まる表現を明示的に決定するには、まず $\mo$ の元から $\mo_0$ における剰余類へ移る。$P$ の元による乗法は $\mo_0$ のすべてのイデアルについて同型である。したがって $P$ は、各単純左イデアルの自己同型斜体 $K_\nu$ の部分体 $e^{(\nu)}P$ に同型である。$P$ 上で $\ml_\nu$ によって媒介される表現は、まずこの部分体 $e^{(\nu)}P$ 上の表現を考え、ついで同型 $e^{(\nu)}P\simeq P$ によってそれを $P$ に移すことで得られる。体 $K_\nu$ は $P$ 上有限階数をもつので、これは $\ml_\nu$ の自己同型斜体の有限次数部分体における表現である。

\[
        c=c_1+\cdots+c_s
\]
を $\mo_0$ の元の両側成分への分解とする。$\ml_\nu$ によって媒介される表現では、$c_\nu$ の行列だけを考える。ほかの $c_i$ は $\ml_\nu$ を零化し、零行列として表現されるからである。$c_\nu$ を行列単位で
\[
        c_\nu=\sum c_{ik}^{(\nu)}\alpha_{ik}^{(\nu)}
\]
と書くなら、各 $\alpha_{ik}^{(\nu)}$ は、$P$ 上の $K_\nu$ の表現においてそれを表す行列で置き換えられるべきである。

$K_\nu$ は $P$ 上有限階数をもつので、$K_\nu$ の任意の元 $x$ は、$e^{(\nu)}P$ に係数をもつ方程式 $f(x)=0$ を満たす。なぜならその冪が線形従属になるからである。$P$ が代数閉であれば、この方程式は一次因子に分解する。$K_\nu$ は体であるから、$x$ はすでにその一次因子の一つの根であり、したがって $e^{(\nu)}P$ に属する。ゆえに $K_\nu=e^{(\nu)}P$ である。この場合には行列 $(\alpha_{ik}^{(\nu)})$ そのものが既約表現である。

これを §19 の終わりと合わせると Burnside の定理を得る。

\emph{$P$ が代数閉であり、環 $\mo$ と可換的に結びついているとする。このとき、$P$ 上の任意の $n$ 次既約表現は、ちょうど $n^2$ 個の線形独立な行列を含む。}

ここから直ちに通常の形が従う。$P$ 上の $n$ 次行列の系で乗法について閉じているものは、すべての行列の成分 $x_{ik}$ が満たす $P$ 係数の非自明な斉次線形関係
\[
        \sum \alpha_{ik}x_{ik}=0
\]
が存在しないとき、かつそのときに限り既約である。実際、すべての線形結合を付け加えると、環かつ $P$-加群が得られ、この環をそれ自身の表現とみなすことができる。

絶対乗法子領域は、単位元をもつ両側単純完全可約環に同型であり、したがってその左イデアルの自己同型斜体上の完全行列環に同型である。$P$ の代数閉性により、その自己同型斜体は $P$ に等しくなる。既約表現は単純左イデアルによって生成される。このイデアルの階数を $m$、表現の次数を $n$ とすると、$m=n$ であり、行列環にはちょうど $n^2$ 個の線形独立な元がある。

同様に一般化 Burnside 定理が従う。同じ仮定の下で、$\mathfrak D$ が次数
\[
        n_1,n_2,\ldots,n_s
\]
の互いに非同値な成分に分解された完全可約表現であるなら、$\mathfrak D$ は $P$ に関して階数
\[
        n_1^2+n_2^2+\cdots+n_s^2
\]
をもつ。

$\mo$ が任意の体上の根基をもたない超複素系であるなら、$\mo$ は完全可約で単位元をもつ。§§17, 18 から、$\mo$ の任意の表現は完全可約である。

逆に、$P$ と可換的に結びついた環 $\mo$ の任意の完全可約表現は、絶対乗法子領域として根基をもたない超複素系をもつ。絶対乗法子領域は、$P$ 上の行列からなる環かつ $P$-加群に同型であり、したがって超複素系である。その根基の元は任意の既約構成成分で零として表現され、したがって全表現でも零として表現される。ゆえに根基は零である。

超複素系 $\mo$ の\emph{正則表現}とは、単位イデアル $\mo$ 自身によって媒介される表現である。群環では、とくに群元を基底として選ぶ。$\mo/\mc$ のすべての左イデアルは $\mo$ の組成因子として現れるので、§16 のように組成列によって正則表現を簡約すれば、すべての既約表現はすでに正則表現の対角上に現れる。

表現は、基底元 $u_1,\ldots,u_h$ を表す行列が分かれば知られる。系が群環であり、したがって $u_i$ が群元であるなら、群の任意の準同型的な行列表現は群環の表現でもある。したがって群を表現する問題は、超複素系を表現する問題の特殊な場合である。

\subsection*{§ 21. 基礎体の拡大。中心の表現}

\[
        \mo=a_1P+\cdots+a_hP
\]
を超複素系とし、$\Omega$ を $P$ の拡大体とする。基底元に対する古い乗法規則を用いて
\[
        \mo\Omega=a_1\Omega+\cdots+a_h\Omega
\]
を作ることができる。\textsuperscript{18a)} このとき $\mo\Omega$ は $\Omega$ に関してふたたび超複素的である。

$\mo$ の任意の表現は $\mo\Omega$ の表現を与える。すべての元の表現は、基底元 $a_i$ の表現が分かれば知られるからである。\textsuperscript{19)}

{\footnotesize\noindent 18a) より正確には、古い乗法規則をもつ新しい記号 $a_i$ を導入する。すると $a_1\Omega+\cdots+a_h\Omega$ は $\mo$ に同型な部分環を含むので、その後で新しい記号を古い $a_i$ と同一視してよい。\par
19) もちろん、ここでも $P$-加群準同型である表現だけを考える。すなわち $a\sim A$ から $\rho\in P$ について $a\rho\sim A\rho$ が従うべきであり、$\Omega$ についても対応して同様である。\par}

イデアルまたは表現加群、および対応する表現は、代数閉体へ移った後も既約にとどまるなら、\emph{絶対既約}と呼ばれる。

$\mo\Omega$ が根基をもたないなら、$\mo$ も根基をもたない。$\mo$ の冪零イデアル $\mc$ は $\mo\Omega$ の拡張イデアル $\mc\Omega$ を与えるからである。逆は成り立たないことが下で分かる。

$\mo\Omega$ の中心は、拡張された中心 $\mZ\Omega$ に等しい。$\mo$ が完全可約であるなら、$\mZ$ は体の直和である：
\[
        \mZ=\mZ_1+\cdots+\mZ_s.
\]
$\Omega$ へ移っても $\mo\Omega$ が完全可約にとどまるなら（すなわち冪零イデアルが加わらないなら）、新しい中心は
\[
        \mZ\Omega=\mZ_1\Omega+\cdots+\mZ_s\Omega
\]
と再び体の直和に分解する。$\Omega$ が代数閉であれば、これらの体は $\Omega$ 上階数 $1$ をもち、$\Omega$ に同型である。

中心 $\mZ$ の表現については、次の定理が成り立つ。

代数閉体 $\Omega$ における可換超複素系 $\mZ$ の任意の既約表現は次数 $1$ である。同値に、既約表現は準同型 $\mZ\to\Omega$ と同一である。

証明。$\mZ$ の任意の既約表現は $\mZ\Omega$ の表現を一つ与え、したがって $\mZ\Omega/\mc$ の表現も与える。ただし $\mc$ は $\mZ\Omega$ の根基である。商 $\mZ\Omega/\mc$ は根基をもたない可換系であり、§14 の終わりにより体の直和である。$\Omega$ は代数閉なので、これらは次数 $1$ であり、すべての既約表現を生成する（§20）。

同時に、次数 $1$ の同値表現は必ず等しくなるので、相異なる準同型 $\mZ\to\Omega$ の数は $\mZ\Omega/\mc$ の階数に等しい。

特別な場合として、$\mZ$ が $P$ 上の体であれば、次の定理が得られる。$\mZ$ が完全可約、すなわち根基をもたないことは、$\mZ$ が $P$ の第一種拡大であることと同値である。体では準同型は同型である。その数は $\mZ\Omega/\mc$ の階数であり、$\mc=0$ のときちょうど体の次数に等しい。

\[
        \mZ=\mZ_1+\cdots+\mZ_s
\]
が完全可約であるなら、各表現において個々の体 $\mZ_i$ も準同型的に表現される。既約表現では、それらの体の一つが同型的に写され、他は零で写される。

根基をもたない系について、これらの事実を超複素系に適用すると、まず次を得る。

$\mo\Omega$、したがって $\mo$ も根基をもたない系であるなら、$\mo$ の任意の絶対既約表現において、中心元 $z$ は対角行列 $E\xi$ によって表現され、$z\mapsto \xi$ は $\Omega$ における $\mZ$ の一次表現である。こうして得られる $\mo$ の絶対既約表現類と $\mZ$ のそれとの対応は一対一である。したがって、これらの表現類の数は中心の階数に等しい。\textsuperscript{20)}

{\footnotesize\noindent 20) 対応する定理は、代数閉体 $\Omega$ における表現だけでなく、個々の自己同型斜体における表現についても成り立つ。これは別の所で展開される予定である。\par}

実際、両側分解不能イデアルへの分解
\[
        \mo\Omega=\ma_1+\cdots+\ma_s
\]
は、中心の体への分解
\[
        \mZ\Omega=\mZ_1+\cdots+\mZ_s
\]
に一対一に対応する。ここで $\mZ_i$ は $\ma_i$ の中心である。$\mo$ の既約表現では、一つの $\ma_i$ が同型的に写され、他は零で写される。表現類は $\ma_i$ によって一意に決まる。成分 $\ma_i$ は完全行列環によって表現され、そのような完全行列環の中心は対角行列 $E\xi$ からなる。対応 $z\mapsto \xi$ は $\mZ$ の準同型であり、一つの $\mZ_i$ は同型的に写され、他は零で写される。これで主張が証明される。

根基のない $\mo$ が根基のない $\mo\Omega$ を与えるのはいつかという問題については、$\mZ$ が $P$ の第二種拡大である成分 $\mZ_i$ をもてば、$\mo\Omega$ は確かに根基をもつことが従う。\textsuperscript{21)} この場合には $\mZ_i\Omega$ がすでに根基をもつからである。

{\footnotesize\noindent 21) $\mZ$ が第一種の成分だけをもつなら、$\mo\Omega$ は根基をもたない。このことも後で証明される。標数零については、I. Schur がすでに同値な定理、すなわち $P$ 上の既約表現は係数領域の任意の拡大の下で完全可約にとどまる、という定理を証明している。I. Schur, Beiträge zur Theorie der Gruppen linearer Substitutionen, \emph{Trans. Amer. Math. Soc.} 15 (1909), p. 159 を参照。\par}

\subsection*{§ 22. アーベル群への応用}

\[
        G=G_1\times G_2\times\cdots\times G_r
\]
を、位数 $h_i$ の巡回群 $G_i$ に分解された有限アーベル群とする。その元は
\[
        a=a_1^{\alpha_1}\cdots a_r^{\alpha_r}
\]
である。$P$ を、標数が群の位数
\[
        h=h_1h_2\cdots h_r
\]
を割らない体とする。群環 $\mZ$ はすべての和
\[
        \sum_{\alpha_1,\ldots,\alpha_r} A_{\alpha_1\ldots\alpha_r}
        a_1^{\alpha_1}\cdots a_r^{\alpha_r}
        \qquad (A\in P)
\]
からなり、$x_i\mapsto a_i$ によって多項式環 $P[x_1,\ldots,x_r]$ の準同型像である。したがって
\[
        \mZ\simeq P[x_1,\ldots,x_r]/(x_1^{h_1}-1,\ldots,x_r^{h_r}-1).
\]
拡大体 $\Omega$ では、多項式 $x_i^{h_i}-1$ は相異なる因子 $x_i-\xi_i$ に分解する。したがって定義イデアルは、相異なる極大イデアル
\[
        (x_1-\xi_1^{(\alpha_1)},\ldots,x_r-\xi_r^{(\alpha_r)})
\]
の積、または交わりである。ただしこれは各零点 $(\xi_1^{(\alpha_1)},\ldots,\xi_r^{(\alpha_r)})$ に一つずつ対応する。この交わりに対して直和分解
\[
        \mZ\Omega=\mZ_1+\cdots+\mZ_h
\]
が対応し、各成分は $\Omega$ に同型であって、表現
\[
        a_i\mapsto \xi_i^{(\alpha_i)},
        \qquad\hbox{またはより一般に}\qquad
        a\mapsto \chi^{(\alpha)}(a)
\]
を与える。これらの表現が指標である。その数は、一般理論が要求する通り、$h_1\cdots h_r=h$ である。

\subsection*{§ 23. 超複素系の行列式}

\[
        \mo=a_1P+\cdots+a_hP
\]
を超複素系とする。$P$ に $h$ 個の不定元 $x_1,\ldots,x_h$ を付け加え、
\[
        \mo^*=a_1P(x)+\cdots+a_hP(x)
\]
を作る。ただし基底元 $a_i$ には古い乗法規則を用いる。$x_i$ は $a_i$ と可換とし、これによって $\mo^*$ における計算規則が定まる。

$\mo^*$ の中には $\mo$ の「一般元」
\[
        w=a_1x_1+\cdots+a_hx_h
\]
がある。表現において $a_i\mapsto A_i$ なら、$w$ には
\[
        W=A_1x_1+\cdots+A_hx_h
\]
が対応する。行列 $W$ をその表現に属する系行列という。$a_i$ が群をなし、$\mo$ が群環であるなら、それは群行列である。正則表現では正則系行列が得られる。

$W$ の成分は $x_i$ の線形形式である。したがって系行列式 $|W|$ は、次数 $n$ の表現に対して次数 $n$ である。とくに正則系行列式は次数 $h$ をもつ。

系行列式は同値な表現へ移っても変わらない。なぜなら
\[
        |PWP^{-1}|=|P|\,|W|\,|P^{-1}|=|W|
\]
だからである。基底 $(a_1,\ldots,a_h)$ から新しい基底 $(b_1,\ldots,b_h)$ へ移り、$w=\sum a_ix_i$ から $w=\sum b_jy_j$ へ移るとき、$W$ の新しい元、したがって新しい行列式も、古いものから変数の正則置換
\[
        x_i=\sum_j \rho_{ij}y_j
\]
によって得られる。

表現加群の組成列が与えられているなら、適切な基底の選択に対して行列 $W$ は下三角ブロック形
\[
        W=\begin{pmatrix}
        W_1&0&\cdots&0\\
        *&W_2&\cdots&0\\
        \vdots&\vdots&\ddots&\vdots\\
        *&*&\cdots&W_s
        \end{pmatrix}
\]
をもつ。任意の表現の行列式 $|W_i|$ の中には、$\mo$ の正則表現、あるいは根基を $\mc$ とする $\mo/\mc$ の正則表現に現れる因子以外は現れない。

$P$ が代数閉であるなら、既約表現に属する行列式 $|W_i|$ は $x_i$ の素関数であり、非同値な表現は異なる素因子を与える。

証明。$\mo$ のすべての既約表現は $\mo/\mc$ の表現でもあるから、根基をもたない環 $\mo/\mc$ に制限してよい。その中で行列単位 $c_{ik}^{(\nu)}$ を基底に取ると、一般元は
\[
        w=\sum_{\nu,i,k} c_{ik}^{(\nu)}x_{ik}^{(\nu)}
\]
である。既約表現の行列は
\[
        W_\nu=(x_{ik}^{(\nu)})
\]
である。関数 $|W_\nu|=|x_{ik}^{(\nu)}|$ は既約であることが知られており、明らかに互いに異なる。

$|W_i|$ を計算するには、$\mo/\mc$ の正則系行列を因数分解すればよい。各素因子は、対応する既約表現の次数だけ現れる。対応するイデアルが組成列の中にその回数だけ現れるからである。

$\mo$ の正則系行列から出発することもできるが、その場合、各既約因子はより多く、すなわち対応する単純イデアルが左イデアルの中で組成因子として現れる回数だけ得られる。この正則表現では $\mo$ は左イデアルとして見られていた。$\mo$ を右イデアルとして見るなら、第二の正則系行列、Frobenius の「反対形行列」が現れる。これは同じ既約因子、すなわちすべての既約表現の系行列式を含むが、指数は異なり得る。

可換系の系行列式は一次因子に分解する。すべての既約表現が次数 $1$ だからである。これらの一次因子は既約表現そのものであり、アーベル群の場合には指標を与える。この事実は、Dedekind が非アーベル群の群行列式を研究する際の出発点であった。

\subsection*{§ 24. トレースと指標}

超複素系 $\mo$ の表現 $\mo\sim\mathfrak D$ において、元 $a$ に行列 $A$ が対応するとき、
\[
        \operatorname{Sp}_D(a)=\operatorname{Sp} A
\]
とおく。トレースは線形関数である：
\[
        \operatorname{Sp}(c+d)=\operatorname{Sp}(c)+\operatorname{Sp}(d),
        \qquad
        \operatorname{Sp}(c\alpha)=\alpha\operatorname{Sp}(c).
\]
同値な表現は同じトレースをもつ。可約表現において、トレースは組成因子によって媒介される表現におけるトレースの和である。実際、
\[
        \operatorname{Sp}\begin{pmatrix}
        A_{11}&0&\cdots&0\\
        *&A_{22}&\cdots&0\\
        \vdots&\vdots&\ddots&\vdots\\
        *&*&\cdots&A_{ss}
        \end{pmatrix}
        =\operatorname{Sp}(A_{11})+\operatorname{Sp}(A_{22})+\cdots+\operatorname{Sp}(A_{ss}).
\]

\emph{主トレース}とは正則表現におけるトレースである。\emph{被約トレース}とは、相異なる既約表現におけるトレースの和である。

$c$ が極大冪零イデアル $\mc$ の元であるなら、任意の表現について $\operatorname{Sp}(c)=0$ である。

証明。$\mo/\mc$ の単純左イデアルによる表現だけを考えれば十分である。任意の他の表現では、これらの組成因子だけが現れ、トレースは加法的に構成されるからである。しかし $c$ は $\mo/\mc$ のすべての元を零化する。したがって各 $c\in\mc$ には零行列が割り当てられ、$\operatorname{Sp}(c)=0$ となる。

$\mo$ が両側単純で、$P$ が代数閉であり、したがって行列環
\[
        \mo=\sum c_{ik}P
\]
であるなら、$\mo$ の既約表現において、元
\[
        a=\sum c_{ik}\alpha_{ik}
\]
には行列 $(\alpha_{ik})$ が割り当てられる。したがって
\[
        \operatorname{Sp}_{\rm red}(a)=\sum_i\alpha_{ii},
        \qquad
        \operatorname{Sp}_{\rm pr}(a)=n\,\operatorname{Sp}_{\rm red}(a)=n\sum_i\alpha_{ii}.
\]
とくに
\[
        \operatorname{Sp}_{\rm red}(c_{ik})=0\quad(i\ne k),\qquad
        \operatorname{Sp}_{\rm red}(c_{ii})=e,
\]
一方で
\[
        \operatorname{Sp}_{\rm pr}(c_{ik})=0\quad(i\ne k),\qquad
        \operatorname{Sp}_{\rm pr}(c_{ii})=n e.
\]
$P$ から代数閉体 $\Omega$ へ移っても主トレースは変わらない。同じ基底から計算できるからである。これは被約トレースについては成り立たない。

根基をもたない系 $\mo$ の元 $a$ の、絶対既約表現におけるトレースを指標と呼び、$\chi(a)$ と書く。表現を指定する場合には $\chi^{(\nu)}(a)$ と書く。

次数 $n_\nu$ の既約表現において、中心元は §21 により対角行列 $E\theta^{(\nu)}$ として表される。ここで $z\mapsto\theta^{(\nu)}(z)$ は中心の次数 $1$ の表現、または中心から $\Omega$ への準同型である。したがってトレースの値は $n_\nu\theta^{(\nu)}(z)$ である。ゆえに中心の準同型 $\theta$ は指標 $\chi$ と
\[
        \chi^{(\nu)}(z)=n_\nu\theta^{(\nu)}(z).              \tag{1}
\]
によって結びついている。可換の場合には $n_\nu=1$ であり、指標自身が準同型を与える。

以下では常にそう仮定するように、$\Omega$ の標数が零であれば、(1) を $n_\nu$ で割ることができる：
\[
        \theta^{(\nu)}(z)=\frac{\chi^{(\nu)}(z)}{n_\nu}.
\]
$\theta^{(\nu)}(z)$ の準同型性は
\[
        \frac{\chi^{(\nu)}(z+z')}{n_\nu}
        =\frac{\chi^{(\nu)}(z)}{n_\nu}
         +\frac{\chi^{(\nu)}(z')}{n_\nu},
\]
および
\[
        \frac{\chi^{(\nu)}(zz')}{n_\nu}
        =\frac{\chi^{(\nu)}(z)}{n_\nu}\,
         \frac{\chi^{(\nu)}(z')}{n_\nu}
\]
で表される。

表現類は、行列のトレースだけですでに一意に定まる。すべての行列のトレースを知るには、もちろん、与えられた表現における基底元のトレースを知れば十分である。

証明。表現加群 $R$ は、各既約表現加群 $\mM_\nu$ がそこに直和因子として何回現れるかを知れば分かる。この数を $\mu_\nu$ とすると、表現における $e^{(\nu)}$ のトレースは $\mu_\nu n_\nu$ である。ただし $n_\nu$ は既約表現の次数である。したがって
\[
        \mu_\nu=\frac{\operatorname{Sp}(e^{(\nu)})}{n_\nu}.
\]

\subsection*{§ 25. 判別式}

\[
        \mo=a_1P+\cdots+a_hP
\]
を超複素系とする。\emph{判別行列}とは、成分が主トレース $\operatorname{Sp}(a_i a_j)$ である行列である。\emph{被約判別行列}とは、代数閉体の中で作られた、被約トレースに関する対応する行列である。この行列の行列式を、それぞれ判別式、被約判別式という。

判別式は基礎体を拡大しても同じである。別の基底に移ると、変換行列式の平方が掛かる。

証明。
\[
        (b_1,\ldots,b_h)=(a_1,\ldots,a_h)P
\]
とする。すると
\[
        \bigl(\operatorname{Sp}(a_i a_j)\bigr)=\bigl(\operatorname{Sp}(a_i b_j)\bigr)P.       \tag{1}
\]
同様に、$P^t$ が転置行列を表すなら
\[
        \bigl(\operatorname{Sp}(a_i b_j)\bigr)=P^t\bigl(\operatorname{Sp}(b_i b_j)\bigr),      \tag{1a}
\]
であり、(1) と (1a) から
\[
        \bigl(\operatorname{Sp}(a_i a_j)\bigr)=P^t\bigl(\operatorname{Sp}(b_i b_j)\bigr)P.
\]
行列式に移れば主張を得る。

したがって行列式は $P$ からの平方因子を除いてしか定まらないが、それが零であるか非零であるかは不変な性質である。(1) からはまた、非零因子を除けば、判別式を二つの異なる基底から、すなわち $|\operatorname{Sp}(a_i b_j)|$ として計算できることも従う。

$\mo$ が冪零イデアル $\mc$ をもてば、判別式は零になる。

証明。$\mo$ の基底元を
\[
        c_1,\ldots,c_r,d_1,\ldots,d_s
\]
の形に選び、$c_1,\ldots,c_r$ が $\mc$ の基底をなすようにする。すると判別行列は、行列式に必要な上ブロックの意味で
\[
        \begin{pmatrix}
        \operatorname{Sp}(c_i c_j)&\operatorname{Sp}(c_i d_j)\\
        \operatorname{Sp}(d_i c_j)&\operatorname{Sp}(d_i d_j)
        \end{pmatrix}
        =
        \begin{pmatrix}
        0&0\\
        0&\operatorname{Sp}(d_i d_j)
        \end{pmatrix}
\]
の形をもつ。したがって行列式は零である。被約判別式についても同じである。

その結果、代数閉体へ移った後に冪零イデアルが現れる場合にも、判別式は零になる。

\[
        \mo=\ma_1+\cdots+\ma_s
\]
を両側直和とし、$M,M_1,\ldots,M_s$ を環 $\mo,\ma_1,\ldots,\ma_s$ の判別行列とする。このとき $M$ はブロック対角であり、
\[
        M=\begin{pmatrix}
        M_1&0&\cdots&0\\
        0&M_2&\cdots&0\\
        \vdots&\vdots&\ddots&\vdots\\
        0&0&\cdots&M_s
        \end{pmatrix},
\]
したがって
\[
        |M|=|M_1|\,|M_2|\cdots |M_s|.
\]
実際、$\mo$ の基底を $\ma_1,\ldots,\ma_s$ の基底から作って選ぶ。$a\in\ma_i$ なら
\[
        \operatorname{Sp}_\mo(a)=\operatorname{Sp}_{\ma_i}(a),
\]
である。ただしどちらも主トレースを意味し、前者は環 $\mo$ におけるもの、後者は環 $\ma_i$ におけるものである。さらに、因子が異なる両側成分に属すれば $a_i a a_j=0$ である。したがって混合トレースブロックは零である。同じ議論が被約判別式にも適用される。

行列環 $\sum c_{ik}P$ の判別式を作るには、二つの異なる基底を選ぶ。すなわち $c_{ik}$ を一度は第一添字で、もう一度は第二添字で順序づける。原典ページに現れる積表は次の通りである：
\[
\begin{array}{c|ccc}
& \overbrace{c_{11}\;\cdots\;c_{1n}}^{r_1}
& \overbrace{c_{21}\;\cdots\;c_{2n}}^{r_2}
& \cdots \\ \hline
\left.\begin{array}{c}c_{11}\\[-1mm]\vdots\\[-1mm]c_{n1}\end{array}\right\}t_1
& (c_{ik}) & 0 & 0 \\[1.3em]
\left.\begin{array}{c}c_{12}\\[-1mm]\vdots\\[-1mm]c_{n2}\end{array}\right\}t_2
& 0 & (c_{ik}) & 0 \\[1.3em]
\vdots & 0 & 0 & (c_{ik})
\end{array}
\]
これから
\[
        D=\left|\begin{array}{cccc}
        \operatorname{Sp}(c_{11})&0&\cdots&0\\
        0&\operatorname{Sp}(c_{22})&\cdots&0\\
        \vdots&\vdots&\ddots&\vdots\\
        0&0&\cdots&\operatorname{Sp}(c_{nn})
        \end{array}\right|^{n}
        =\begin{cases}
        e, & \text{被約トレースの場合},\\
        n^{n^2}e, & \text{主トレースの場合}.
        \end{cases}
\]
すなわち
\[
        D_{\rm red}=e
        \quad\hbox{被約トレースについて},
        \qquad
        D=n^{n^2}e
        \quad\hbox{主トレースについて}.
\]

したがって、代数閉拡大体上で $\mo$ が次数 $n_1,\ldots,n_s$ の行列環に分解するなら、その判別式は
\[
        D=n_1^{n_1^2}n_2^{n_2^2}\cdots n_s^{n_s^2}e,
\]
またその被約判別式は
\[
        D_{\rm red}=e.
\]
したがって、根基をもたない系について被約判別式は常に非零である。主判別式は、どの $n_i$ も体の標数で割り切れないとき、かつそのときに限り非零である。ゆえに次を得る。

\emph{被約判別式が非零であることは、体 $P$ の代数閉包へ移った後に根基をもたない系であるための必要十分条件である。}

\emph{標数零では、判別式が非零であることは、体 $P$ の代数閉包へ移った後に根基が現れないための必要十分条件である。}\textsuperscript{22)}

{\footnotesize\noindent 22) 可換系については、E. Noether, \emph{Diskriminantensatz für Ordnungen ...}, J. reine angew. Math. 157 (1927), pp. 82--104, §§4--6 を参照。そこの証明方法は同じであるが、細部ではより複雑である。基底から別の基底への移行（同論文 §4, 4）は、ここでこの節の冒頭に与えたものに置き換えるべきである。そこに現れる行列式は零になるからである。\par}

\subsection*{§ 26. 群環の位置づけ}

$a_1,\ldots,a_h$ を有限群の元とし、標数が $h$ を割らない体上で群環を作る。このとき判別式 $D\ne0$ であり、したがって群環は根基をもたない環である。

証明。まず、以後 $\operatorname{Sp}$ を主トレースとすることにして、
\[
        \operatorname{Sp}(e)=he,
        \qquad
        \operatorname{Sp}(a_i)=0\quad (a_i\ne e)
\]
を示す。正則表現では
\[
        e\mapsto\begin{pmatrix}e&0&\cdots&0\\0&e&\cdots&0\\\vdots&\vdots&\ddots&\vdots\\0&0&\cdots&e\end{pmatrix},
        \qquad \operatorname{Sp}(e)=he
\]
である。$a_i\ne e$ に対して、積 $a_i a_1,\ldots,a_i a_h$ は $a_1,\ldots,a_h$ の固定点のない置換である。したがって、これらの積を $a_1,\ldots,a_h$ によって表す行列は主対角上に零だけをもつので、$\operatorname{Sp}(a_i)=0$ である。

いま基底 $a_1,\ldots,a_h$ と $a_1^{-1},\ldots,a_h^{-1}$ を取る。行列 $\bigl(\operatorname{Sp}(a_i a_j^{-1})\bigr)$ は対角行列で、対角成分は $he$ である。したがって
\[
        D=h^h e\ne0.
\]

群元 $a$ の類とは、群環の中で作られる和
\[
        K_a=\sum_s s^{-1}as,                              \tag{1}
\]
をいう。ただし和は相異なる共役元 $s^{-1}as$ だけにわたる。類 $K_a$ は中心元である。すべての群元と可換だからである。これらは中心を生成する。実際、群環の元 $\sum_i \alpha_i a_i$ がすべての $s$ と可換であるなら
\[
        \sum_i\alpha_i a_i=s^{-1}\biggl(\sum_i\alpha_i a_i\biggr)s
        =\sum_i\alpha_i(s^{-1}a_i s),
\]
であるから、係数は各共役類上で一定である。

したがって中心の階数は共役な群元の類の数に等しい。ゆえに絶対既約表現の数もこの類数に等しい。

行列 $A$ と $S^{-1}AS$ は同じトレースをもつ。したがって群元 $a$ と $s^{-1}as$ は同じトレースをもつ。(1) の両辺でトレースを取ると
\[
        \chi(K_a)=h_a\chi(a),
\]
を得る。ここで $h_a$ は類 $K_a$ に含まれる元の数である。言葉で言えば、群元の指標は、その類の指標を類に含まれる元の数で割ったものに等しい。

\begin{center}
（1928年8月12日受理。）
\end{center}


\end{document}
